% ============================================================
%  module_01_genetics.tex  —  IBO Biology Platform
%  Module 1: Genetics and Molecular Biology
%
%  STANDALONE: compile this file directly on Overleaf.
%  Upload to Overleaf together with:
%    preamble.tex  +  the four .jpg images in this folder.
%  Set compiler: pdfLaTeX, root: module_01_genetics.tex
% ============================================================

\documentclass[12pt, a4paper, oneside]{report}
% ============================================================
%  preamble.tex  —  IBO Biology Platform
%  Minimal academic style. pdfLaTeX compatible.
% ============================================================

% --- Page geometry ---
\usepackage[a4paper,
            top=2.5cm, bottom=2.8cm,
            left=3cm, right=2.5cm]{geometry}

% --- Typography ---
\usepackage[T1]{fontenc}
\usepackage[utf8]{inputenc}
\usepackage{lmodern}
\usepackage{microtype}
\usepackage{setspace}
\setstretch{1.12}

% --- Mathematics ---
\usepackage{amsmath}
\usepackage{amssymb}
\usepackage{amsfonts}
\usepackage{mathtools}
\usepackage{bm}
\usepackage{siunitx}
\sisetup{per-mode=symbol}

% --- Chemistry ---
\usepackage[version=4]{mhchem}

% --- Graphics ---
\usepackage{graphicx}
\usepackage{xcolor}
\usepackage{tikz}
\usetikzlibrary{calc}

% --- Floats and captions ---
\usepackage{float}
\usepackage{caption}
\captionsetup{font=small, labelfont=bf, labelsep=period, width=0.92\textwidth}

% --- Tables ---
\usepackage{booktabs}
\usepackage{tabularx}
\usepackage{array}

% --- Lists ---
\usepackage{enumitem}
\setlist[enumerate,1]{label=\textbf{(\alph*)}, itemsep=2pt}
\setlist[enumerate,2]{label=\roman*.,  itemsep=1pt}
\setlist[itemize]{itemsep=2pt}

% --- Framed environments (for solutions, notes) ---
\usepackage{mdframed}

% --- Colour palette: one accent, otherwise near-black ---
\definecolor{accent}{HTML}{1A2E4A}    % dark navy
\definecolor{probback}{HTML}{F5F6F7}  % near-white box background
\definecolor{stargold}{HTML}{A07820}  % muted gold
\definecolor{rulecol}{HTML}{AAAAAA}   % light rule colour
\definecolor{codebg}{HTML}{F8F8F8}
\definecolor{codeframe}{HTML}{CCCCCC}
\definecolor{codekw}{HTML}{00397A}
\definecolor{codestr}{HTML}{B83200}
\definecolor{codecomment}{HTML}{4A7A4A}
\definecolor{codenumber}{HTML}{999999}

% --- Hyperlinks ---
\usepackage{hyperref}
\hypersetup{
    colorlinks = true,
    linkcolor  = accent,
    citecolor  = accent,
    urlcolor   = accent,
    pdftitle   = {IBO Biology Platform},
}

% --- Code listings ---
\usepackage{listings}
\usepackage{lstautogobble}

\lstdefinestyle{python}{
    language          = Python,
    backgroundcolor   = \color{codebg},
    basicstyle        = \ttfamily\small,
    keywordstyle      = \color{codekw}\bfseries,
    stringstyle       = \color{codestr},
    commentstyle      = \color{codecomment}\itshape,
    numberstyle       = \tiny\color{codenumber},
    numbers           = left,
    stepnumber        = 1,
    numbersep         = 8pt,
    frame             = single,
    framesep          = 4pt,
    rulecolor         = \color{codeframe},
    breaklines        = true,
    showstringspaces  = false,
    tabsize           = 4,
    captionpos        = b,
    autogobble        = true,
    morekeywords      = {as, with, None, True, False,
                         np, plt, odeint, linspace, exp, log, pi, array},
}
\lstset{style=python}

% --- tcolorbox (problems only) ---
\usepackage{tcolorbox}
\tcbuselibrary{breakable, skins}

% Problem box: clean, numbered, minimal
\newtcolorbox[auto counter, number within=section]{problem}[1]{%
    breakable,
    enhanced,
    colback    = probback,
    colframe   = accent,
    boxrule    = 0.5pt,
    arc        = 1pt,
    fonttitle  = \bfseries\sffamily\small,
    coltitle   = white,
    colbacktitle = accent,
    title      = {Problem~\thetcbcounter.\enspace #1},
    attach boxed title to top left={yshift=-1.5mm, xshift=3mm},
    boxed title style={arc=1pt, boxrule=0pt},
    top        = 5pt, bottom = 5pt,
    left       = 6pt, right  = 6pt,
}

% Solution: simple left-rule indent, no background
\newenvironment{solution}{%
    \begin{mdframed}[
        linewidth    = 1.5pt,
        linecolor    = accent!40,
        topline      = false,
        rightline    = false,
        bottomline   = false,
        leftmargin   = 4pt,
        innerleftmargin = 10pt,
        innertopmargin  = 4pt,
        innerbottommargin = 4pt,
        backgroundcolor = white,
    ]%
    \noindent\textit{Solution.}\enspace\ignorespaces
}{%
    \end{mdframed}%
}

% Historical note: small italic, thin left rule
\newenvironment{historicalnote}[1]{%
    \begin{mdframed}[
        linewidth    = 2pt,
        linecolor    = rulecol,
        topline      = false,
        rightline    = false,
        bottomline   = false,
        leftmargin   = 4pt,
        innerleftmargin = 10pt,
        innertopmargin  = 3pt,
        innerbottommargin = 3pt,
        backgroundcolor = white,
    ]%
    {\small\sffamily\bfseries Historical note: #1}\\[3pt]%
    \small\itshape\ignorespaces
}{%
    \end{mdframed}%
}

% Definition: bold label, thin left rule
\newenvironment{definition}[1]{%
    \begin{mdframed}[
        linewidth    = 1.5pt,
        linecolor    = accent!60,
        topline      = false,
        rightline    = false,
        bottomline   = false,
        leftmargin   = 4pt,
        innerleftmargin = 10pt,
        innertopmargin  = 3pt,
        innerbottommargin = 3pt,
        backgroundcolor = white,
    ]%
    \noindent\textbf{Definition (#1).}\enspace\ignorespaces
}{%
    \end{mdframed}%
}

% Key remark: just bold inline label, no box
\newenvironment{keyinsight}{%
    \par\medskip\noindent\textbf{Remark.}\enspace\ignorespaces
}{\par\medskip}

% Model limitations: just bold inline label
\newenvironment{limitation}{%
    \par\medskip\noindent\textbf{Limitations of this model.}\enspace\ignorespaces
}{\par\medskip}

% Summary: simple ruled block
\newenvironment{summary}{%
    \begin{mdframed}[
        linewidth = 0.5pt, linecolor = accent!50,
        leftmargin = 0pt, innerleftmargin = 8pt,
        innertopmargin = 6pt, innerbottommargin = 6pt,
        backgroundcolor = probback,
    ]%
    \noindent\textbf{Summary.}\enspace\ignorespaces
}{%
    \end{mdframed}%
}

% Code lab: just a listing caption prefix
\newenvironment{codelab}[1]{%
    \par\medskip\noindent\textbf{Code.}\enspace\textit{#1}\par\smallskip
}{\par}

% ============================================================
%  Difficulty stars — plain \ifnum, no loops
% ============================================================
\usepackage{pifont}

\newcommand{\starfilled}{{\color{stargold}\ding{72}}}
\newcommand{\starempty} {{\color{stargold!40}\ding{73}}}

\newcommand{\stars}[1]{%
  \ifnum#1>0\starfilled\else\starempty\fi
  \ifnum#1>1\starfilled\else\starempty\fi
  \ifnum#1>2\starfilled\else\starempty\fi
  \ifnum#1>3\starfilled\else\starempty\fi
  \ifnum#1>4\starfilled\else\starempty\fi
}

% ============================================================
%  Discipline tags — plain \fcolorbox, safe everywhere
% ============================================================
\newcommand{\disciplinetag}[1]{%
  {\setlength{\fboxsep}{2pt}\setlength{\fboxrule}{0.4pt}%
   \fcolorbox{accent!50}{accent!8}{%
     \textcolor{accent}{\sffamily\tiny\bfseries\vphantom{Ag}#1}}}%
}

% ============================================================
%  Section / chapter headings — clean, minimal
% ============================================================
\usepackage{titlesec}

\titleformat{\chapter}[block]
  {\normalfont\LARGE\bfseries\sffamily\color{accent}}
  {Module~\thechapter.}{0.6em}{}[\vspace{2pt}\rule{\textwidth}{0.4pt}]

\titleformat{\section}
  {\normalfont\large\bfseries\sffamily\color{accent}}
  {\thesection}{0.7em}{}

\titleformat{\subsection}
  {\normalfont\normalsize\bfseries\sffamily}
  {\thesubsection}{0.6em}{}

\titleformat{\subsubsection}
  {\normalfont\normalsize\itshape}
  {}{0em}{}

% ============================================================
%  Header / footer — minimal
% ============================================================
\usepackage{fancyhdr}
\pagestyle{fancy}
\fancyhf{}
\fancyhead[RO,LE]{\small\sffamily\thepage}
\fancyhead[LO]{\small\sffamily\nouppercase{\rightmark}}
\fancyhead[RE]{\small\sffamily\nouppercase{\leftmark}}
\renewcommand{\headrulewidth}{0.3pt}
\renewcommand{\footrulewidth}{0pt}

% ============================================================
%  Maths macros
% ============================================================
\newcommand{\kB}{k_{\mathrm{B}}}
\newcommand{\NA}{N_{\mathrm{A}}}
\newcommand{\dG}{\Delta G}
\newcommand{\ddG}{\Delta\Delta G}
\newcommand{\dH}{\Delta H}
\newcommand{\dS}{\Delta S}
\newcommand{\Kd}{K_{\mathrm{d}}}
\newcommand{\Ka}{K_{\mathrm{a}}}
\newcommand{\Keq}{K_{\mathrm{eq}}}
\newcommand{\vmax}{V_{\mathrm{max}}}
\newcommand{\Km}{K_{\mathrm{m}}}
\newcommand{\conc}[1]{[\mathrm{#1}]}
\newcommand{\od}[2]{\frac{\mathrm{d}#1}{\mathrm{d}#2}}
\newcommand{\pd}[2]{\frac{\partial #1}{\partial #2}}
\newcommand{\e}{\mathrm{e}}
\newcommand{\abs}[1]{\left|#1\right|}

% ============================================================
%  Text abbreviations
% ============================================================
\newcommand{\Ecoli}{\textit{E.~coli}}
\newcommand{\invitro}{\textit{in vitro}}
\newcommand{\invivo}{\textit{in vivo}}
\newcommand{\etal}{\textit{et al.}}

% End of preamble.tex


\begin{document}

% ---------- Title block ----------
\begin{titlepage}
  \centering
  \vspace*{4cm}
  {\sffamily\fontsize{13}{16}\selectfont\color{accent} IBO Biology Platform}\\[0.6em]
  {\rule{0.4\textwidth}{0.4pt}}\\[0.6em]
  {\sffamily\fontsize{22}{28}\selectfont\bfseries Module 1}\\[0.4em]
  {\sffamily\fontsize{16}{20}\selectfont\bfseries\color{accent} Genetics and Molecular Biology}\\[1.2em]
  {\small\color{gray}
    DNA structure \textbullet\ Replication fidelity \textbullet\
    Transcription kinetics \textbullet\ Translation \textbullet\
    Gene regulation \textbullet\ DNA repair}\\[0.6em]
  {\rule{0.4\textwidth}{0.4pt}}\\[0.5em]
  {\small\color{gray} \today}
\end{titlepage}

\tableofcontents
\newpage

% ============================================================
%  Chapter begins here
% ============================================================

\chapter{Genetics and Molecular Biology}
\label{ch:genetics}

\begin{tcolorbox}[colback=probback, colframe=accent, boxrule=0.5pt, arc=1pt,
  fonttitle=\bfseries\sffamily\small, coltitle=white, colbacktitle=accent,
  title={Module Overview}, top=4pt, bottom=4pt, left=6pt, right=6pt]
  \textbf{Topics:} DNA structure and information content;
  replication machinery and fidelity;
  transcription initiation kinetics;
  translation and the genetic code;
  gene regulation (\textit{lac} operon);
  DNA repair.\\[3pt]
  \textbf{Prerequisites:} Basic thermodynamics and equilibrium chemistry;
  introductory calculus.\\[3pt]
  \textbf{Problem set:} 10 problems, \stars{2}--\stars{5}.
\end{tcolorbox}

% ============================================================
\section{The Physical Structure of DNA}
% ============================================================

\subsection{Information storage requirements}

A cell must store information sufficient to specify all its proteins and
regulatory elements, and must copy that information faithfully at every
division. Three physical requirements follow: stable information storage,
high-fidelity copying, and regulated readout. DNA addresses all three.

\begin{historicalnote}{Watson, Crick, Franklin and the Double Helix (1953)}
  The structure of DNA was determined in 1953 at the Cavendish Laboratory,
  Cambridge. James Watson and Francis Crick built their model using
  Rosalind Franklin's X-ray diffraction data (``Photo 51''), obtained
  at King's College London, and Erwin Chargaff's empirical rules
  ($[\text{A}]=[\text{T}]$; $[\text{G}]=[\text{C}]$). Franklin's data
  directly revealed the helical parameters: pitch 3.4~nm, repeat 0.34~nm
  per base pair, diameter 2.0~nm. Watson and Crick received the Nobel
  Prize in Physiology or Medicine in 1962, jointly with Maurice Wilkins.
  Franklin died in 1958 and was not eligible for posthumous nomination.
\end{historicalnote}

\subsection{Chemical structure and directionality}

DNA is a polymer of four deoxyribonucleotides: adenine (A), thymine (T),
guanine (G), and cytosine (C). The phosphodiester backbone connects the
3$'$ hydroxyl of one deoxyribose to the 5$'$ phosphate of the next,
creating a directional chain. The two strands run \textbf{antiparallel}
and are held together by Watson--Crick hydrogen bonds:
\[
  \text{A--T}: \; 2 \text{ H-bonds}, \qquad \text{G--C}: \; 3 \text{ H-bonds}.
\]
Because the strands are complementary and antiparallel, each strand
serves as a complete template for synthesising the other.


\subsection{Thermodynamics of duplex stability}

The stability of the double helix does not arise primarily from hydrogen
bonds between bases. Hydrogen bonds to water are also favourable, so
the enthalpic gain from base pairing is partially offset by the loss of
hydrogen bonds to solvent. The dominant stabilising force is
\textbf{base stacking}: London dispersion interactions between the
stacked $\pi$-electron systems of adjacent aromatic rings, combined
with the hydrophobic burial of the ring surfaces from water.

The empirical melting temperature $T_m$ of a short DNA duplex in solution:
\begin{equation}
  T_m \approx 81.5 + 16.6\log_{10}[\text{Na}^+] + 0.41(\%\text{GC})
             - \frac{675}{n} \quad (^\circ\text{C}),
  \label{eq:tm}
\end{equation}
where $[\text{Na}^+]$ is the salt concentration (M), $\%\text{GC}$ is
the fraction of G$+$C base pairs, and $n$ is the duplex length (bp).
Salt screens repulsion between the negatively charged backbones (raises
$T_m$); G--C pairs are more stable than A--T pairs; longer duplexes
accumulate more stacking interactions.

\subsection{Information content of DNA}

The four nucleotides form a four-letter alphabet. The information content
per symbol drawn uniformly from an alphabet of $k$ symbols is
$I = \log_2 k$ bits (Shannon, 1948). With $k = 4$:
\[
  I = \log_2 4 = 2 \text{ bits per base pair.}
\]
The human genome ($\approx 3.2 \times 10^9$~bp) therefore stores
$\approx 800$~MB of information. The storage density of DNA
($\sim10^{21}$ bits~cm$^{-3}$) exceeds that of current electronic
memory by many orders of magnitude.

% ============================================================
\section{DNA Replication: Machinery and Fidelity}
% ============================================================

\subsection{The Replication Fork}

\begin{figure}[H]
  \centering
  \includegraphics[width=0.92\textwidth]{dna_replication_fork_1786906792977.jpg}
  \caption{The \Ecoli\ replication fork. DnaB helicase (yellow hexamer) unwinds
           the parental duplex. SSBs coat the lagging-strand template.
           DNA Pol~III holoenzyme synthesises the leading strand continuously
           and the lagging strand in Okazaki fragments, each primed by
           DnaG primase.}
  \label{fig:repfork}
\end{figure}

Replication begins at \textit{oriC} in \Ecoli. The initiator protein DnaA
melts a region of \textit{oriC} and loads the helicase DnaB, which forms
a hexameric ring encircling the lagging-strand template and unwinds the
helix at $\approx 1000$~bp~s$^{-1}$.

The key components and their roles are summarised in
Table~\ref{tab:repfork}.

\begin{table}[H]
  \centering
  \caption{Major protein components of the \Ecoli\ replication fork.}
  \label{tab:repfork}
  \small
  \begin{tabularx}{\textwidth}{lXl}
    \toprule
    \textbf{Protein} & \textbf{Function} & \textbf{Rate / size} \\
    \midrule
    DnaB helicase  & Unwinds parental duplex (5$'\to$3$'$ on lagging template)
                   & $\sim$1000~bp~s$^{-1}$ \\
    SSBs           & Coat ssDNA, prevent re-annealing, remove secondary structure
                   & Tetrameric \\
    DnaG primase   & Synthesises 10--12~nt RNA primers \textit{de novo}
                   & Every $\sim$1--2~kb on lagging strand \\
    Pol~III holoenzyme & Replicative polymerase (5$'\to$3$'$ synthesis only)
                   & $\sim$1000~nt~s$^{-1}$; processivity $>500$~kb \\
    $\beta$ clamp  & Ring-shaped processivity clamp encircling DNA
                   & Dimer \\
    Pol~I          & Removes RNA primers (5$'\to$3$'$ exonuclease) and fills gap
                   & \\
    DNA Ligase     & Seals nicks between Okazaki fragments (uses NAD$^+$)
                   & \\
    \bottomrule
  \end{tabularx}
\end{table}

\subsubsection{The Directionality Problem}

All known DNA polymerases synthesise exclusively 5$'\to$3$'$. Since the
template strands are antiparallel, only the \textbf{leading strand} can
be synthesised continuously in the direction of fork movement. The
\textbf{lagging strand} is synthesised in discontinuous 1--2~kb segments
called \emph{Okazaki fragments} (100--200~nt in eukaryotes), each
requiring a new RNA primer. This asymmetry is a direct physical
consequence of chemical directionality.

\begin{historicalnote}{Reiji Okazaki and the Discovery of Discontinuous Replication (1968)}
  In 1968, Reiji Okazaki at Nagoya University labelled newly synthesised
  DNA with short pulses of $^3$H-thymidine and analysed the products by
  alkaline sucrose gradient centrifugation. He observed a population of
  short (1--2~kb) fragments, which matured into high-molecular-weight DNA
  with longer labelling times. This directly demonstrated discontinuous
  synthesis on the lagging strand. The short intermediates bear his name.
  Okazaki died in 1975 from radiation-induced cancer, a consequence of
  atomic bomb exposure in Hiroshima.
\end{historicalnote}

\subsection{Replication Fidelity: Three Tiers of Error Correction}

The spontaneous error rate of nucleotide insertion in solution is
$\sim10^{-1}$--$10^{-2}$ per base (due to tautomeric base shifts). The
final biological error rate is $\sim10^{-9}$--$10^{-10}$ per base — a
$10^{10}$-fold improvement achieved through three successive mechanisms.

\subsubsection*{Tier 1 — Intrinsic Selectivity of Pol III: $\varepsilon_1 \approx 10^{-5}$}

The active site of Pol~III selects the correct dNTP through an
induced-fit mechanism. Correct base-pair geometry triggers closure of
the ``fingers'' domain, properly positioning the 3$'$-OH and the
incoming $\alpha$-phosphate for catalysis. Mismatched nucleotides
prevent full domain closure, reducing the rate of phosphodiester bond
formation by $10^4$--$10^5$-fold.

\subsubsection*{Tier 2 — 3$'\to$5$'$ proofreading: $\sim$100-fold reduction}

The $\varepsilon$ subunit of the Pol~III core carries a 3$'\to$5$'$
exonuclease. After each nucleotide addition, there is kinetic competition
between chain extension and transfer to the exonuclease site. A mismatched
3$'$ end is extended $\sim10^3$-fold more slowly than a correct end,
shifting the terminus preferentially to the exonuclease, which removes it.

\subsubsection*{Tier 3 — Methyl-directed Mismatch Repair (MMR): $\sim$100-fold reduction}

\begin{figure}[H]
  \centering
  \includegraphics[width=0.92\textwidth]{dna_mismatch_repair_1786906827550.jpg}
  \caption{The \Ecoli\ methyl-directed mismatch repair system.
           MutS recognises the geometric distortion of the mismatch.
           MutH identifies the daughter strand by the absence of
           Dam methylation at GATC sites. The error-containing segment
           is excised by UvrD and exonucleases, then resynthesised by Pol~III.}
  \label{fig:mmr}
\end{figure}

After replication, any remaining mismatches are identified by MutS. The
critical problem of strand discrimination — which strand contains the
error? — is solved by \textbf{Dam methylation}. Dam methylase modifies
adenine within GATC sequences. Immediately after replication, the parental
strand is methylated, but the daughter strand is not yet methylated
(the hemimethylated window persists for $\sim$1--4 minutes). MutH
specifically nicks the \emph{unmethylated} strand at hemimethylated
GATC sites, thereby targeting excision to the newly synthesised strand
that carries the error.

\begin{table}[H]
  \centering
  \caption{Three-tier error correction in \Ecoli.}
  \label{tab:fidelity}
  \begin{tabular}{llr}
    \toprule
    \textbf{Mechanism} & \textbf{Error rate after} & \textbf{Reduction} \\
    \midrule
    None (uncorrected)        & $\sim10^{-1}$       & ---        \\
    Pol III selectivity       & $\sim10^{-5}$       & $10^4$     \\
    3$'\to$5$'$ proofreading  & $\sim10^{-7}$       & $10^2$     \\
    Mismatch repair           & $\sim10^{-9}$       & $10^2$     \\
    \bottomrule
  \end{tabular}
\end{table}

% ============================================================
\section{Transcription: From DNA to RNA}
% ============================================================

\subsection{Promoter recognition}

\begin{figure}[H]
  \centering
  \includegraphics[width=0.92\textwidth]{transcription_initiation_1786906804692.jpg}
  \caption{Transcription initiation in \Ecoli. The $\sigma^{70}$ factor
           contacts the $-35$ (TTGACA) and $-10$ (TATAAT) promoter elements,
           enabling holoenzyme ($\alpha_2\beta\beta'\omega\sigma$) to form the
           closed complex. Local strand melting from $-11$ to $+3$ creates the
           open complex. RNA synthesis begins at $+1$.}
  \label{fig:transcription}
\end{figure}

The $\sigma^{70}$ factor contacts two conserved sequence elements on the
nontemplate (coding) strand:
\begin{itemize}
  \item \textbf{$-35$ element} (consensus \texttt{TTGACA}): bound by the
        $\sigma_4$ domain.
  \item \textbf{$-10$ element} (consensus \texttt{TATAAT}): bound by the
        $\sigma_2$ domain. The AT-rich composition is functionally important
        — A--T base pairs melt more easily than G--C pairs (2 vs.\ 3
        hydrogen bonds), lowering the activation barrier for the open
        complex.
  \item \textbf{Spacer} of $17 \pm 1$~bp between $-35$ and $-10$. Deviation
        from 17~bp rotates the relative orientation of the two $\sigma$
        domains, reducing simultaneous contact and weakening promoter
        strength.
\end{itemize}

\subsection{Kinetics of initiation}

Initiation proceeds through a sequence of states:
\begin{equation}
  R + P \underset{k_{-1}}{\overset{k_1}{\rightleftharpoons}} RP_c
        \underset{k_{-2}}{\overset{k_2}{\rightleftharpoons}} RP_o
        \xrightarrow{k_3} \text{initial transcribing complex} \to \cdots
  \label{eq:txn_kinetics}
\end{equation}
where $RP_c$ is the closed complex (RNAP bound, DNA duplex intact) and
$RP_o$ is the open complex (transcription bubble from $-11$ to $+3$
melted). The isomerization $RP_c \to RP_o$ is typically rate-limiting
and thermally activated. Under the pre-equilibrium assumption
($k_1, k_{-1} \gg k_2$):

\begin{definition}{Observed Rate of Open Complex Formation}
  \[
    k_\text{obs} = \frac{k_2 \, [\text{RNAP}]}{[\text{RNAP}] + K_d^c}
    \quad\text{where}\quad K_d^c = \frac{k_{-1}}{k_1}.
  \]
  As $[\text{RNAP}] \to \infty$, $k_\text{obs} \to k_2$, the intrinsic
  isomerization rate.
\end{definition}

\subsection{Termination}

Two mechanisms terminate transcription in bacteria:
\begin{enumerate}
  \item \textbf{Intrinsic (Rho-independent):} A GC-rich stem-loop forms in
        the nascent RNA near the 3$'$ end, followed by a run of U residues.
        The rU--dA hybrid is thermodynamically weak, causing spontaneous
        RNA release.
  \item \textbf{Rho-dependent:} The Rho hexameric helicase loads onto the
        5$'$ end of the nascent RNA and translocates 5$'\to$3$'$,
        powered by ATP hydrolysis. When RNAP pauses at a terminator,
        Rho catches up and unwinds the RNA--DNA hybrid.
\end{enumerate}

% ============================================================
\section{Translation: Decoding the Transcript}
% ============================================================

\subsection{The Ribosome}

The 70S bacterial ribosome consists of the 30S small subunit (16S rRNA +
$\sim$21 proteins) and the 50S large subunit (23S rRNA + 5S rRNA +
$\sim$31 proteins). The small subunit decodes mRNA; the large subunit
catalyses peptide bond formation. The peptidyl transferase centre is
composed entirely of RNA — the ribosome is a ribozyme.

\begin{historicalnote}{Ada Yonath and the Ribosome Structure (2009 Nobel Prize)}
  Determining the atomic structure of the ribosome required decades of
  methodological innovation. Ada Yonath (Weizmann Institute) pioneered
  cryo-cooling of ribosome crystals to reduce radiation damage during
  X-ray diffraction. Her group published the first ribosome crystal
  structures in the 1980s. High-resolution structures ($<3$~\AA) were
  ultimately obtained by Yonath, Thomas Steitz, and Venkatraman
  Ramakrishnan, who shared the 2009 Nobel Prize in Chemistry. These
  structures revealed that the active site contains no protein side
  chains within 18~\AA\ of the peptidyl transfer reaction — definitively
  confirming the ribosome as a ribozyme.
\end{historicalnote}

\subsection{The Genetic Code and Its Information Content}

The genetic code uses triplet codons ($4^3 = 64$ possible) to encode
20 amino acids and 3 stop signals. The information content per codon:
\[
  I_\text{codon} = \log_2 64 = 6 \text{ bits}.
\]
The minimum required to distinguish 20 amino acids (assuming equal
frequencies):
\[
  I_\text{min} = \log_2 20 \approx 4.32 \text{ bits}.
\]
The \textbf{redundancy} of the code is $6 - 4.32 = 1.68$~bits~codon$^{-1}$,
corresponding to $\approx 3.2$-fold degeneracy on average (consistent
with 64 codons / 20 amino acids). Crucially, synonymous codons nearly
always differ only at the \textbf{third (wobble) position}, concentrating
mutational buffering at the position where tRNA wobble pairing is most
tolerant of non-Watson--Crick contacts.

\subsection{Energetic Cost of Translation}

Each peptide bond costs approximately 4 ATP equivalents:
\begin{itemize}
  \item Aminoacyl-tRNA synthesis: ATP $\to$ AMP $+$ PP$_\text{i}$
        (equivalent to 2 ATP in high-energy bonds).
  \item EF-Tu GTP hydrolysis for aminoacyl-tRNA delivery to A site.
  \item EF-G GTP hydrolysis for translocation.
\end{itemize}
A 300-amino-acid protein therefore costs $\sim$1200 ATP molecules.

% ============================================================
\section{Gene Regulation: The Thermodynamics of the \textit{lac} Operon}
% ============================================================

\begin{figure}[H]
  \centering
  \includegraphics[width=0.95\textwidth]{lac_operon_regulation_1786906815973.jpg}
  \caption{\textit{E.~coli} \textit{lac} operon regulation. \textbf{Left:}
           Repressed state (no lactose). The lac repressor tetramer bridges
           operators O1, O2, and O3, looping the DNA and blocking RNAP.
           \textbf{Right:} Induced state (lactose/IPTG present). Allolactose
           binds the repressor and reduces its DNA affinity $\sim10^4$-fold,
           allowing RNAP to transcribe \textit{lacZYA}.}
  \label{fig:lacoperon}
\end{figure}

\subsection{Quantitative repressor binding}

\begin{definition}{Fractional Occupancy}
  For the equilibrium $R + O \rightleftharpoons RO$ with dissociation
  constant $\Kd = \conc{R}\conc{O}/\conc{RO}$, the fraction of operators
  occupied by repressor is:
  \[
    \theta = \frac{\conc{R}}{\conc{R} + \Kd}.
  \]
\end{definition}

For the \textit{lac} repressor: $\Kd^\text{op} \approx 10^{-13}$~M;
intracellular $\conc{R} \approx 10~\text{nM} = 10^{-8}$~M.

\[
  \theta_\text{rep} = \frac{10^{-8}}{10^{-8} + 10^{-13}}
                    = \frac{1}{1 + 10^{-5}} \approx 0.99999.
\]

Addition of IPTG reduces the DNA-binding affinity by $\sim10^4$-fold
($\Kd^\text{ind} \approx 10^{-9}$~M):
\[
  \theta_\text{ind} = \frac{10^{-8}}{10^{-8} + 10^{-9}}
                    = \frac{10}{11} \approx 0.909.
\]

The rate of transcript initiation is proportional to the unoccupied
promoter fraction $(1-\theta)$. The theoretical fold-induction is:
\[
  \text{Fold-induction} = \frac{1 - \theta_\text{ind}}{1 - \theta_\text{rep}}
  = \frac{0.091}{10^{-5}} \approx 9{,}100.
\]
The experimentally measured value is $\sim$1000-fold, reflecting protein
dynamics, incomplete IPTG saturation, and in vivo conditions deviating
from in vitro $\Kd$ measurements.

\subsection{Catabolite repression}

When glucose is absent, cAMP levels rise. cAMP binds the Catabolite
Activator Protein (CAP/CRP), which contacts the $\alpha$CTD of RNAP and
increases the open complex formation rate $\sim$20-fold.
The \textit{lac} operon is therefore expressed at high levels only when
both conditions hold:
\[
  \text{High expression} \Longleftrightarrow
  (\text{glucose absent}) \;\text{AND}\; (\text{lactose present}).
\]

\subsection{Cooperative binding and the Hill equation}

When multiple ligand-binding events are positively cooperative, the
fractional occupancy follows the Hill equation:
\begin{equation}
  \theta = \frac{\conc{L}^n}{\conc{L}^n + \Kd^n},
  \label{eq:hill}
\end{equation}
where $n$ is the Hill coefficient ($n=1$ gives the Michaelis--Menten
hyperbola; $n > 1$ indicates positive cooperativity, sharpening the
transition around $\Kd$). The lac repressor tetramer, bridging operators
O1, O2, and O3 via DNA looping, gives an effective $n \approx 2$.

% ============================================================
\section{Mutation and DNA Repair}
% ============================================================

Types of mutations include:
\begin{enumerate}
  \item \textbf{Point mutations} — transitions (purine $\leftrightarrow$
        purine, or pyrimidine $\leftrightarrow$ pyrimidine; more frequent)
        and transversions (purine $\leftrightarrow$ pyrimidine; less frequent).
  \item \textbf{Indels} — insertions or deletions; cause frameshifts if
        not multiples of 3.
  \item \textbf{Silent (synonymous)} — no amino acid change.
  \item \textbf{Missense} — altered amino acid.
  \item \textbf{Nonsense} — premature stop codon.
\end{enumerate}

Repair pathways include:
\begin{itemize}
  \item \textbf{Base Excision Repair (BER):} Removes small base lesions
        (deaminated cytosine $\to$ uracil, oxidised guanine). DNA glycosylase
        cleaves the glycosidic bond, AP endonuclease nicks the backbone,
        Pol $\beta$ fills, and ligase seals.
  \item \textbf{Nucleotide Excision Repair (NER):} Removes bulky lesions
        (UV dimers, cisplatin adducts). Excises a $\sim$12--25~nt segment
        around the lesion (bacteria) or $\sim$25--30~nt (eukaryotes).
  \item \textbf{Mismatch Repair (MMR):} Described in Section~1.2.
  \item \textbf{Double-Strand Break Repair:} Via homologous recombination
        (high fidelity) or non-homologous end joining (error-prone).
\end{itemize}

% ============================================================
\section{Code Lab: Gene Expression Dynamics}
% ============================================================

\begin{codelab}{\textit{lac} operon induction and protein accumulation}
  Two simulations: (a)~protein dilution after translational block by
  chloramphenicol; (b)~\textit{lac} operon induction from repressed state
  to new steady state.
\end{codelab}

\begin{lstlisting}[language=Python, caption={Gene expression dynamics (run in Python 3 with NumPy, SciPy, and Matplotlib installed).}, label={lst:geneexpr}]
import numpy as np
import matplotlib.pyplot as plt
from scipy.integrate import odeint

# -------------------------------------------------------
# Parameters
# -------------------------------------------------------
k_trans  = 10.0           # translation rate (proteins / mRNA / min)
delta_M  = np.log(2) / 2  # mRNA degradation (1/min; half-life = 2 min)
delta_P  = np.log(2) / 30 # protein dilution (1/min; doubling time = 30 min)
M_ss     = 10.0           # steady-state mRNA (molecules / cell)

# Steady-state protein copy number
P_ss = k_trans * M_ss / delta_P
print(f"Steady-state mRNA:    {M_ss:.0f}  molecules/cell")
print(f"Steady-state protein: {P_ss:.0f} molecules/cell")
print(f"mRNA half-life:  {np.log(2)/delta_M:.1f} min")
print(f"Protein half-life: {np.log(2)/delta_P:.1f} min  (= cell doubling time)")

# -------------------------------------------------------
# Part 1: Protein decay after chloramphenicol (t = 0)
# mRNA degrades exponentially; protein only lost by dilution
# -------------------------------------------------------
def model_cam(P, t, k_trans, M0, delta_M, delta_P):
    M = M0 * np.exp(-delta_M * t)  # mRNA decays
    return k_trans * M - delta_P * P

t1    = np.linspace(0, 60, 1000)
sol1  = odeint(model_cam, [P_ss], t1,
               args=(k_trans, M_ss, delta_M, delta_P))

# -------------------------------------------------------
# Part 2: lac operon induction dynamics
# IPTG added at t = 30 min: effective transcription rate
# jumps from 1% to 100% of maximum
# -------------------------------------------------------
k_mRNA = delta_M * M_ss   # transcription rate sustaining M_ss at SS

def model_lac(state, t, k_mRNA, delta_M, k_trans, delta_P, t_iptg):
    M, P  = state
    k_tx  = k_mRNA if t >= t_iptg else 0.01 * k_mRNA  # repressed -> induced
    dMdt  = k_tx - delta_M * M
    dPdt  = k_trans * M - delta_P * P
    return [dMdt, dPdt]

t2     = np.linspace(0, 120, 3000)
init   = [0.1, 43.0]   # nearly fully repressed initial state
sol2   = odeint(model_lac, init, t2,
                args=(k_mRNA, delta_M, k_trans, delta_P, 30.0))

# -------------------------------------------------------
# Plotting
# -------------------------------------------------------
fig, axes = plt.subplots(1, 2, figsize=(13, 5))

# Panel A: chloramphenicol response
ax = axes[0]
ax.axhline(P_ss, color='gray', ls=':', lw=1.2,
           label=f'$P_{{ss}} = {P_ss:.0f}$ molecules')
ax.plot(t1, sol1[:, 0], 'k-', lw=2, label='Protein (Cm added at $t=0$)')
ax.axvline(np.log(2)/delta_P, color='red', ls='--', lw=1.2,
           label=f'Protein $t_{{1/2}} = {np.log(2)/delta_P:.1f}$ min')
ax.set_xlabel('Time (min)', fontsize=12)
ax.set_ylabel('Protein molecules per cell', fontsize=12)
ax.set_title('(A) Protein decay after translational block', fontsize=11)
ax.legend(fontsize=9)
ax.grid(True, alpha=0.3)

# Panel B: lac induction
ax = axes[1]
ax.axvline(30, color='orange', ls='--', lw=1.5, label='IPTG added')
ax.plot(t2, sol2[:, 0], 'b--', lw=1.8, label='mRNA (molecules/cell)')
ax.plot(t2, sol2[:, 1], 'r-',  lw=2.2, label='Protein (molecules/cell)')
ax.set_xlabel('Time (min)', fontsize=12)
ax.set_ylabel('Molecules per cell', fontsize=12)
ax.set_title('(B) $lac$ operon induction dynamics', fontsize=11)
ax.legend(fontsize=9)
ax.grid(True, alpha=0.3)

plt.suptitle('Gene Expression Dynamics', fontsize=13, fontweight='bold')
plt.tight_layout()
plt.savefig('gene_expression_dynamics.pdf', dpi=150)
plt.show()
\end{lstlisting}

\textbf{Results.}
\begin{itemize}
  \item Panel~(A): mRNA decays with $t_{1/2} = 2$~min after translational
        block, but protein decays with $t_{1/2} = 30$~min (the doubling
        time). Protein is only removed by dilution at division.
  \item Panel~(B): mRNA reaches its new steady state within
        $\sim$5~min of IPTG addition. Protein follows more slowly,
        equilibrating after $\sim$3--4 doubling times ($\approx$90--120~min).
\end{itemize}

% ============================================================
%  PROBLEM SET
% ============================================================
\newpage
\section{Problem Set: Module 1}

\noindent\textbf{Instructions.} Attempt each problem before reading the
solution. Show all working. Where numbers are requested, carry appropriate
significant figures and include units throughout.

% ------------------------------------------------------------
\begin{problem}{\stars{2}\enspace\disciplinetag{Replication Fidelity}\enspace\disciplinetag{Quantitative}\enspace Mutation rate arithmetic}

  A DNA polymerase has an intrinsic error rate of
  $\varepsilon_1 = 10^{-5}$ per base incorporated.
  The 3$'\to$5$'$ proofreading exonuclease reduces this by a factor of
  $f_1 = 10^2$. Subsequent mismatch repair reduces it by a further factor
  of $f_2 = 10^2$.

  \begin{enumerate}
    \item Calculate the final per-base error rate $\varepsilon_\text{final}$.
    \item The \Ecoli\ genome is $G = 4.6 \times 10^6$~bp. What is the
          expected number of mutations per complete genome replication?
    \item \Ecoli\ divides every $T = 20$~min under optimal conditions.
          How many mutations accumulate in a single lineage over one year?
    \item A culture contains $N = 10^{10}$ cells. Calculate the expected
          number of cells carrying a pre-existing mutation at a specific
          single nucleotide. Comment on antibiotic resistance.
  \end{enumerate}
\end{problem}

\begin{solution}
  \textbf{(a)} The three correction mechanisms act sequentially and
  independently:
  \[
    \varepsilon_\text{final}
    = \varepsilon_1 \times \frac{1}{f_1} \times \frac{1}{f_2}
    = 10^{-5} \times 10^{-2} \times 10^{-2}
    = \boxed{10^{-9} \text{ per base.}}
  \]

  \textbf{(b)}
  \[
    \mu = \varepsilon_\text{final} \times G
        = 10^{-9} \times 4.6 \times 10^6
        = 4.6 \times 10^{-3} \text{ mutations per replication.}
  \]
  On average, one complete genome replication in \Ecoli\ produces fewer
  than one mutation.

  \textbf{(c)} Divisions per year:
  \[
    n_\text{div} = \frac{365 \times 24 \times 60}{20} = 26{,}280.
  \]
  \[
    \mu_\text{year} = 4.6 \times 10^{-3} \times 26{,}280 \approx 121
    \text{ mutations per year.}
  \]

  \textbf{(d)} Expected cells harbouring a specific mutation:
  \[
    E = N \times \varepsilon_\text{final}
      = 10^{10} \times 10^{-9} = 10 \text{ cells.}
  \]
  Ten pre-existing mutant cells are expected before any antibiotic is
  applied. Antibiotic treatment does not create resistance — it \emph{selects}
  for it. This is why insufficient antibiotic doses (selecting but not
  sterilising) are particularly dangerous.
\end{solution}

% ------------------------------------------------------------
\begin{problem}{\stars{3}\enspace\disciplinetag{Gene Regulation}\enspace\disciplinetag{Thermodynamics}\enspace Lac repressor binding}

  The \textit{lac} repressor binds operator O1 with
  $\Kd = 10^{-13}$~M. The intracellular repressor concentration
  is $\conc{R} = 10$~nM.

  \begin{enumerate}
    \item Derive the fractional occupancy equation
          $\theta = \conc{R}/(\conc{R} + \Kd)$ from equilibrium
          thermodynamics. State all assumptions.
    \item Calculate $\theta$ for the repressed state.
    \item Adding 1~mM IPTG shifts the effective $\Kd$ by a factor of $10^4$.
          Calculate $\Kd^\text{ind}$ and $\theta_\text{ind}$.
    \item The rate of transcript initiation is proportional to
          $(1-\theta)$. Calculate the theoretical fold-induction.
    \item The experimentally measured fold-induction is $\approx$1000-fold.
          Is this consistent with (d)? Identify possible sources of
          discrepancy.
  \end{enumerate}
\end{problem}

\begin{solution}
  \textbf{(a)} The binding equilibrium is:
  \[
    R + O \underset{K_d}{\rightleftharpoons} RO, \qquad
    \Kd = \frac{\conc{R}\conc{O}}{\conc{RO}}.
  \]
  Defining $\theta = \conc{RO}/(\conc{O}+\conc{RO})$ and substituting
  $\conc{O} = \Kd\conc{RO}/\conc{R}$:
  \[
    \theta
    = \frac{\conc{RO}}{\Kd\conc{RO}/\conc{R} + \conc{RO}}
    = \frac{1}{\Kd/\conc{R} + 1}
    = \frac{\conc{R}}{\conc{R} + \Kd}. \qquad\square
  \]
  Assumptions: (i) equilibrium applies; (ii) $\conc{R} \gg \conc{O}$
  so free repressor concentration $\approx$ total repressor
  (valid: $\sim$10 repressor copies per cell vs.\ 1 operator).

  \textbf{(b)}
  \[
    \theta_\text{rep}
    = \frac{10^{-8}}{10^{-8} + 10^{-13}}
    = \frac{1}{1 + 10^{-5}}
    \approx 0.99999.
  \]

  \textbf{(c)} $\Kd^\text{ind} = 10^4 \times 10^{-13} = 10^{-9}$~M.
  \[
    \theta_\text{ind}
    = \frac{10^{-8}}{10^{-8} + 10^{-9}}
    = \frac{10}{11} \approx 0.909.
  \]

  \textbf{(d)}
  \[
    \text{Fold-induction}
    = \frac{1 - 0.909}{1 - 0.99999}
    = \frac{0.091}{10^{-5}}
    = \boxed{9{,}100}.
  \]

  \textbf{(e)} Our prediction ($\sim$9100) exceeds the measured $\sim$1000
  by $\approx$10-fold. The model overestimates because: (i) \invitro\
  $\Kd$ values do not perfectly reflect \invivo\ conditions (macromolecular
  crowding, binding of DNA-associated proteins); (ii) IPTG does not fully
  saturate all repressor at the concentrations used; (iii) at full
  induction, the maximum promoter strength limits transcription
  independently of repressor; (iv) the model assumes perfectly sharp
  induction but repressor unbinding is kinetically finite.
  The order-of-magnitude agreement validates the thermodynamic framework.
\end{solution}

% ------------------------------------------------------------
\begin{problem}{\stars{3}\enspace\disciplinetag{Translation Kinetics}\enspace\disciplinetag{ODE Modelling}\enspace mRNA and protein dynamics}

  An mRNA reaches a steady-state concentration of 10 molecules/cell
  (half-life $t_{1/2}^M = 2$~min). Each mRNA is translated at
  $k_\text{trans} = 10$~proteins~min$^{-1}$. Proteins are lost by
  cell division ($T_d = 30$~min; negligible intrinsic degradation).

  \begin{enumerate}
    \item Write the ODE for protein copy number $P(t)$ under steady-state
          mRNA. Find $P_\text{ss}$.
    \item Chloramphenicol blocks translation at $t = 0$. mRNA decays as
          $M(t) = M_0\e^{-\delta_M t}$. Write and solve the resulting ODE
          for $P(t)$.
    \item How long does it take for protein to fall to $P_\text{ss}/2$?
          Why does this differ from the mRNA half-life?
  \end{enumerate}
\end{problem}

\begin{solution}
  \textbf{(a)} Let $\delta_M = \ln2/2$~min$^{-1}$ and
  $\delta_P = \ln2/T_d = \ln2/30$~min$^{-1}$. At steady-state mRNA:
  \[
    \od{P}{t} = k_\text{trans} M_\text{ss} - \delta_P P.
  \]
  Setting $\od{P}{t} = 0$:
  \[
    P_\text{ss} = \frac{k_\text{trans} M_\text{ss}}{\delta_P}
                = \frac{10 \times 10}{\ln2/30}
                \approx \boxed{4{,}330 \text{ proteins/cell.}}
  \]

  \textbf{(b)} With mRNA decaying, the ODE becomes:
  \[
    \od{P}{t} = k_\text{trans} M_0 \e^{-\delta_M t} - \delta_P P.
  \]
  Multiply both sides by integrating factor $\e^{\delta_P t}$:
  \[
    \od{}{t}\!\left(P\e^{\delta_P t}\right)
    = k_\text{trans} M_0 \e^{(\delta_P - \delta_M)t}.
  \]
  Integrate and apply $P(0) = P_\text{ss}$:
  \[
    P(t) = \frac{k_\text{trans} M_0}{\delta_P - \delta_M}
           \!\left(\e^{-\delta_M t} - \e^{-\delta_P t}\right)
           + P_\text{ss}\e^{-\delta_P t}.
  \]
  Since $\delta_M \gg \delta_P$, the mRNA is exhausted after
  $\sim$3$\times t_{1/2}^M \approx 6$~min. Thereafter:
  \[
    P(t) \approx P_\text{ss}\,\e^{-\delta_P t}.
  \]

  \textbf{(c)} The protein half-life under dilution is
  $t_{1/2}^P = \ln2/\delta_P = 30$~min — equal to the cell
  doubling time. This far exceeds the 2-min mRNA half-life because
  \textbf{protein is only removed by cell division}, not by rapid
  intracellular degradation. mRNA and protein operate on completely
  different timescales, which enables cells to respond to stimuli
  at the mRNA level (fast) while maintaining protein levels stably
  over cell generations.
\end{solution}

% ------------------------------------------------------------
\begin{problem}{\stars{3}\enspace\disciplinetag{Information Theory}\enspace\disciplinetag{Genetic Code}\enspace The information content of the genetic code}

  \begin{enumerate}
    \item Calculate the minimum bits required to encode 20 amino acids.
    \item A triplet code over 4 nucleotides encodes $4^3 = 64$ codons.
          Calculate the code's information content and its redundancy in bits.
    \item Synonymous codons usually differ only at the wobble (third)
          position. Explain why this pattern reduces the mutational burden.
    \item Given a translational error rate of $\varepsilon = 10^{-4}$
          per codon, use the Poisson distribution to calculate the
          probability that a 300-amino-acid protein is synthesised with
          (i) zero errors and (ii) exactly one error.
  \end{enumerate}
\end{problem}

\begin{solution}
  \textbf{(a)} Minimum bits to distinguish 20 equally probable amino acids:
  \[
    I_\text{min} = \log_2 20 \approx 4.32 \text{ bits.}
  \]

  \textbf{(b)} Information per codon:
  $I_\text{codon} = \log_2 64 = 6$~bits.
  Redundancy $= 6 - 4.32 = 1.68$~bits~codon$^{-1}$, corresponding to
  $\approx 3.2$-fold degeneracy on average.

  \textbf{(c)} If synonymous codons differ only at position 3, then a
  point mutation at the wobble position has a high probability of remaining
  \emph{synonymous} (encoding the same amino acid). Mutations at positions
  1 or 2 almost always change the amino acid. The code therefore concentrates
  mutational tolerance at the wobble position, where tRNA wobble pairing
  is tolerant of G--U and other non-Watson--Crick pairs anyway. This is
  widely thought to be a selected feature of the code, reducing the mean
  fitness cost of both replication errors and translational misreading.

  \textbf{(d)} Expected errors per protein: $\lambda = 300 \times 10^{-4} = 0.03$.
  Poisson:
  \begin{align*}
    P(0 \text{ errors}) &= \e^{-0.03} \approx 0.9704, \\
    P(1 \text{ error})  &= 0.03\,\e^{-0.03} \approx 0.0291.
  \end{align*}
  $\approx$97\% of 300-aa proteins are synthesised without error.
\end{solution}

% ------------------------------------------------------------
\begin{problem}{\stars{3}\enspace\disciplinetag{Transcription}\enspace\disciplinetag{Thermodynamics}\enspace Promoter binding and activation}

  A promoter has $\Kd^\text{RNAP} = 10^{-8}$~M for the
  $\sigma^{70}$-holoenzyme. Intracellular holoenzyme concentration
  $\conc{H} = 50$~nM.

  \begin{enumerate}
    \item Calculate the fractional RNAP occupancy at this promoter.
    \item A transcriptional activator stabilises RNAP at the promoter by
          $\ddG = -3.0$~kcal~mol$^{-1}$ at $T = 310$~K.
          Using $R = 1.987$~cal~mol$^{-1}$~K$^{-1}$, calculate the new
          effective $\Kd^\text{act}$.
    \item A repressor ($\Kd^\text{rep} = 5\times10^{-10}$~M,
          $\conc{Rp} = 100$~nM) competes with RNAP for the same site.
          Using a partition function over three states (empty, RNAP-bound,
          repressor-bound), derive and calculate the probability of RNAP
          occupancy.
  \end{enumerate}
\end{problem}

\begin{solution}
  \textbf{(a)}
  \[
    \theta = \frac{5\times10^{-8}}{5\times10^{-8} + 10^{-8}} = \frac{5}{6} \approx 0.833.
  \]

  \textbf{(b)} The free energy change shifts the association constant:
  \[
    \frac{\Kd^\text{act}}{\Kd^\text{basal}} = \e^{\ddG/RT}, \qquad
    RT = 1.987 \times 310 = 616 \text{ cal mol}^{-1}.
  \]
  \[
    \frac{\Kd^\text{act}}{10^{-8}} = \e^{-3000/616} = \e^{-4.87}
    \approx 7.7 \times 10^{-3},
  \]
  \[
    \Kd^\text{act} \approx 7.7 \times 10^{-11} \text{ M} \approx 0.077 \text{ nM.}
  \]
  A 130-fold reduction in $\Kd$ converts a moderately occupied promoter
  ($\theta = 0.833$) to one that is nearly always occupied.

  \textbf{(c)} Define statistical weights for each state:
  \[
    Z = 1 + \frac{\conc{H}}{\Kd^\text{RNAP}} + \frac{\conc{Rp}}{\Kd^\text{rep}}.
  \]
  \[
    \frac{\conc{H}}{\Kd^\text{RNAP}} = \frac{5\times10^{-8}}{10^{-8}} = 5; \qquad
    \frac{\conc{Rp}}{\Kd^\text{rep}} = \frac{10^{-7}}{5\times10^{-10}} = 200.
  \]
  \[
    P_\text{RNAP} = \frac{5}{1 + 5 + 200} = \frac{5}{206} \approx 0.024.
  \]
  RNAP occupies the promoter only 2.4\% of the time — the repressor,
  with a 40-fold greater occupancy weight, almost entirely excludes RNAP.
\end{solution}

% ------------------------------------------------------------
\begin{problem}{\stars{4}\enspace\disciplinetag{Molecular Evolution}\enspace\disciplinetag{Information Theory}\enspace\disciplinetag{Population Genetics}\enspace The Drake rule and the error threshold}

  \begin{enumerate}
    \item The Drake rule states that the spontaneous mutation rate is
          $\mu \approx 0.003$ per genome per replication across
          DNA-based microorganisms. For \Ecoli\ ($G = 4.6\times10^6$~bp),
          calculate the implied per-base error rate and verify consistency
          with Table~\ref{tab:fidelity}.
    \item For poliovirus ($G_\text{RNA} = 7.4\times10^3$~nt, RdRp error
          rate $\approx 10^{-4}$~nt$^{-1}$), calculate the mutations per
          genome replication. Does poliovirus follow the Drake rule?
    \item Define the Eigen error threshold mathematically. Using
          $L = 7.4\times10^3$ and average mutant fitness $w_m = 0.1$,
          calculate the threshold error rate $\mu^*$.
    \item Ribavirin approximately doubles the RdRp error rate. Determine
          whether a 2-fold increase is sufficient to push poliovirus
          above its error threshold.
  \end{enumerate}
\end{problem}

\begin{solution}
  \textbf{(a)}
  \[
    \varepsilon = \mu / G = 0.003 / (4.6\times10^6) \approx 6.5\times10^{-10}
    \text{ per base,}
  \]
  consistent with the MMR-corrected value of $\sim10^{-9}$--$10^{-10}$
  in Table~\ref{tab:fidelity}.

  \textbf{(b)} Poliovirus: $10^{-4} \times 7.4\times10^3 \approx 0.74$
  mutations per replication — 250-fold above the Drake rule. RNA viruses
  exist at the edge of tolerable mutation load; they maintain fitness
  as a quasispecies cloud rather than a single consensus sequence.

  \textbf{(c)} The probability of copying a genome of length $L$ without
  error: $Q = (1-\mu)^L \approx \e^{-\mu L}$.
  At the error threshold, the wild-type cannot maintain informational
  dominance:
  \[
    Q^* \approx w_m \implies \e^{-\mu^* L} = w_m \implies
    \mu^* = -\frac{\ln w_m}{L} = \frac{\ln 10}{7400} \approx 3.1\times10^{-4}
    \text{ per nt.}
  \]

  \textbf{(d)} Current poliovirus error rate: $10^{-4}$~nt$^{-1}$.
  With ribavirin: $2\times10^{-4}$~nt$^{-1}$.
  Threshold: $3.1\times10^{-4}$~nt$^{-1}$.

  \[
    2\times10^{-4} < 3.1\times10^{-4}.
  \]
  A 2-fold increase is \textbf{insufficient}. A $\gtrsim$3.1-fold increase
  would be required. This is consistent with clinical observations that
  ribavirin has modest antiviral efficacy against poliovirus as a
  monotherapy, and explains why combination therapy is typically required
  for lethal mutagenesis strategies.
\end{solution}

% ------------------------------------------------------------
\begin{problem}{\stars{3}\enspace\disciplinetag{DNA Thermodynamics}\enspace Duplex stability and promoter function}

  A 20-bp duplex with alternating AT/GC has an estimated hybridisation
  enthalpy of $\dH = -160$~kJ~mol$^{-1}$ (simplified nearest-neighbour
  calculation).

  \begin{enumerate}
    \item Using $\dS = \dH / 354$~K (empirical approximation) and
          $T_m = \dH / (\dS + R\ln(C_T/4))$ with
          $C_T = 1~\mu$M, calculate $T_m$ in kelvin and $^\circ$C.
    \item The $-10$ element of a bacterial promoter has the consensus
          sequence \texttt{TATAAT} (100\% AT). Explain mechanistically
          why this AT-richness is a functional requirement rather than
          an evolutionary accident.
    \item A mutant promoter has the $-10$ sequence changed to
          \texttt{GCGCGC} (100\% GC). Predict the qualitative effect on
          (i) $T_m$ of the $-10$ region, (ii) open complex formation rate,
          and (iii) transcriptional output.
  \end{enumerate}
\end{problem}

\begin{solution}
  \textbf{(a)}
  \[
    \dS = \frac{-160{,}000}{354} = -452~\text{J mol}^{-1}\text{K}^{-1}.
  \]
  \[
    R\ln(C_T/4) = 8.314\ln(2.5\times10^{-7}) = -126~\text{J mol}^{-1}\text{K}^{-1}.
  \]
  \[
    T_m = \frac{-160{,}000}{-452 + (-126)} = \frac{-160{,}000}{-578}
        \approx 277~\text{K} = 4^\circ\text{C.}
  \]
  (This simplified calculation underestimates real $T_m$ values; accurate
  results require the full SantaLucia 1998 nearest-neighbour parameters.)

  \textbf{(b)} Open complex formation requires RNAP to locally melt
  $\sim$14~bp of DNA from $-11$ to $+3$. This melting requires thermal
  energy to break base-pair hydrogen bonds and disrupt stacking. AT-rich
  sequences have lower melting enthalpy and entropy than GC-rich sequences;
  they are easier to melt at physiological temperatures. The consensus
  \texttt{TATAAT} sequence at the $-10$ element therefore lowers the
  activation barrier for the closed-to-open complex isomerization step,
  increasing the rate of productive transcription initiation. This is
  a thermodynamic necessity — not an arbitrary sequence preference.

  \textbf{(c)} A \texttt{GCGCGC} $-10$ element: (i) $T_m$ of the $-10$
  region increases substantially (100\% GC content raises $T_m$ by
  $\approx$4°C per 10\% GC increase). (ii) The higher melting energy
  raises the activation barrier for open complex formation, greatly
  reducing $k_\text{iso}$. (iii) Transcriptional output is predicted to
  fall sharply — this promoter would be classified as very weak. In practice,
  such a mutation would likely be lethal if the gene is essential, or
  would cause a strong loss-of-expression phenotype.
\end{solution}

% ------------------------------------------------------------
\begin{problem}{\stars{4}\enspace\disciplinetag{Population Genetics}\enspace\disciplinetag{Molecular Evolution}\enspace Kimura's neutral theory and the molecular clock}

  \begin{enumerate}
    \item The rate of neutral substitution $k$ in a population is
          equal to the neutral mutation rate $\mu_n$, independent of
          population size $N$. Derive this result from first principles:
          the rate of new neutral mutations per generation is $2N\mu_n$
          (diploid), and the fixation probability of a neutral allele
          is $1/(2N)$.
    \item Two mammalian species share a gene with 10\% synonymous
          divergence. Using $\mu_n = 2\times10^{-9}$~site$^{-1}$~yr$^{-1}$
          and the Jukes--Cantor correction, estimate their divergence time.
    \item Why must synonymous (silent) substitutions — rather than all
          substitutions — be used as the molecular clock? What systematic
          bias would occur if nonsynonymous substitutions were included?
  \end{enumerate}
\end{problem}

\begin{solution}
  \textbf{(a)} In a population of $N$ diploid individuals, each generation
  produces $2N\mu_n$ new neutral mutations. Each new neutral allele starts
  at frequency $1/(2N)$ and has fixation probability equal to its initial
  frequency (for a neutral allele, $p_\text{fix} = 1/(2N)$ by the
  random walk argument). The rate of neutral substitutions per generation:
  \[
    k = (\text{mutations per generation}) \times p_\text{fix}
      = 2N\mu_n \times \frac{1}{2N} = \mu_n. \qquad\square
  \]
  The population size cancels exactly — a profound result. Large populations
  generate more neutral mutations, but each has a proportionally smaller
  fixation probability.

  \textbf{(b)} Jukes--Cantor correction for multiple hits at the same site:
  \[
    d = -\frac{3}{4}\ln\!\left(1 - \frac{4p}{3}\right), \quad p = 0.10.
  \]
  \[
    d = -0.75\ln(0.867) \approx 0.107.
  \]
  Divergence time (factor of 2 for independent accumulation on both lineages):
  \[
    t = \frac{d}{2\mu_n} = \frac{0.107}{2\times2\times10^{-9}}
      \approx 2.68\times10^7 \text{ yr} \approx 26.8 \text{ Mya.}
  \]

  \textbf{(c)} Synonymous substitutions are approximately neutral (no
  change in protein sequence, so minimal selection) and accumulate at
  rate $\approx\mu_n$. Nonsynonymous substitutions are under selection:
  deleterious ones are eliminated (purifying selection reduces their rate
  below $\mu_n$) and advantageous ones sweep through rapidly. Including
  nonsynonymous substitutions introduces rate heterogeneity that depends
  on the specific functional constraints and evolutionary history of
  the gene, making the clock unreliable. Using all substitutions would
  typically \emph{underestimate} divergence times (purifying selection
  suppresses the apparent substitution rate).
\end{solution}

% ------------------------------------------------------------
\begin{problem}{\stars{3}\enspace\disciplinetag{Transcription Kinetics}\enspace Open complex formation rate}

  For the \textit{lac} promoter: $k_\text{on} = 2\times10^7$~M$^{-1}$s$^{-1}$,
  $k_\text{off} = 0.2$~s$^{-1}$, $k_\text{iso} = 1$~s$^{-1}$,
  $k_{-\text{iso}} = 0.01$~s$^{-1}$, $\conc{H} = 50$~nM.

  \begin{enumerate}
    \item Under the pre-equilibrium assumption
          ($k_\text{on}, k_\text{off} \gg k_\text{iso}$), derive an
          expression for the observed rate $k_\text{obs}$ of open
          complex formation as a function of $\conc{H}$.
    \item Calculate $k_\text{obs}$ and $\tau = 1/k_\text{obs}$ at
          $\conc{H} = 50$~nM.
    \item What is the physical interpretation of the maximum rate
          $k_\text{obs}^\text{max}$ as $\conc{H} \to \infty$?
  \end{enumerate}
\end{problem}

\begin{solution}
  \textbf{(a)} The closed complex occupancy under pre-equilibrium:
  \[
    \theta_c = \frac{\conc{H}}{\conc{H} + K_d^c},
    \quad K_d^c = k_\text{off}/k_\text{on}
               = 0.2/(2\times10^7) = 10^{-8} \text{ M.}
  \]
  The observed isomerization rate:
  \[
    k_\text{obs} = k_\text{iso}\,\theta_c
    = \frac{k_\text{iso}\conc{H}}{\conc{H} + K_d^c}.
  \]

  \textbf{(b)} At $\conc{H} = 5\times10^{-8}$~M:
  \[
    k_\text{obs} = \frac{1 \times 5\times10^{-8}}{5\times10^{-8} + 10^{-8}}
                 = \frac{5}{6} \approx 0.833 \text{ s}^{-1}.
  \]
  \[
    \tau = 1/0.833 \approx 1.2 \text{ s.}
  \]
  Open complex formation takes $\approx$1.2~s under physiological conditions.

  \textbf{(c)} As $\conc{H} \to \infty$, all promoters are saturated with
  closed complexes and $k_\text{obs} \to k_\text{iso} = 1$~s$^{-1}$.
  This is the intrinsic rate of the $RP_c \to RP_o$ conformational change
  — determined by the activation energy of DNA melting, not by RNAP
  concentration. No amount of additional polymerase can accelerate the
  reaction beyond this intrinsic limit.
\end{solution}

% ------------------------------------------------------------
\begin{problem}{\stars{5}\enspace\disciplinetag{Replication}\enspace\disciplinetag{DNA Repair}\enspace\disciplinetag{Cancer Biology}\enspace Mismatch repair, microsatellite instability, and Lynch syndrome}

  Lynch syndrome is caused by heterozygous germline loss-of-function
  mutations in MMR genes (MLH1, MSH2, MSH6, PMS2) and confers
  $\approx$70--80\% lifetime colorectal cancer risk.

  \begin{enumerate}
    \item MMR-deficient cells operate with only two correction tiers.
          Calculate the per-base error rate $\varepsilon_\text{MMR^-}$
          and the expected mutations per human genome replication
          ($G = 3\times10^9$~bp).
    \item Lynch syndrome tumours show microsatellite instability (MSI) —
          length variation at short tandem repeat sequences. Explain
          mechanistically why MMR deficiency specifically elevates
          microsatellite mutation rates rather than uniformly raising
          all mutation types.
    \item Model colorectal carcinogenesis as requiring $m = 5$
          rate-limiting mutations at specific single base positions.
          For a colon with $N = 10^{10}$ epithelial cells undergoing
          $n = 100$ divisions per lifetime, estimate the probability
          that at least one cell has accumulated all 5 mutations, for:
          \begin{enumerate}[label=(\roman*)]
            \item MMR-proficient cells ($\varepsilon = 10^{-9}$~bp$^{-1}$)
            \item MMR-deficient cells ($\varepsilon = 10^{-7}$~bp$^{-1}$)
          \end{enumerate}
    \item Your estimates in (c) are far smaller than observed cancer
          incidences. Identify the most important physical process
          missing from this simple model.
  \end{enumerate}
\end{problem}

\begin{solution}
  \textbf{(a)} Without MMR, only two tiers remain:
  \[
    \varepsilon_{\text{MMR}^-}
    = \varepsilon_1 / f_1
    = 10^{-5} \times 10^{-2}
    = 10^{-7} \text{ per base.}
  \]
  Mutations per replication ($G = 3\times10^9$~bp):
  \[
    \mu_{\text{MMR}^-} = 10^{-7} \times 3\times10^9 = 300 \text{ per replication,}
  \]
  compared to $\mu_\text{normal} = 10^{-9} \times 3\times10^9 = 3$.
  MMR loss produces a 100-fold mutator phenotype.

  \textbf{(b)} Microsatellites (tandem repeat sequences such as CACACACA)
  are particularly prone to \textbf{replication slippage}: during Okazaki
  fragment synthesis on the lagging strand, the newly synthesised strand
  can loop out (or the template can loop out), causing the polymerase to
  re-prime at the wrong repeat unit, inserting or deleting one repeat
  unit. This creates a small insertion/deletion loop (IDL) rather than a
  base substitution mismatch. IDLs are the primary substrate of the
  MSH2--MSH3 mismatch recognition complex. Without functional MMR,
  these IDLs are not corrected, causing length variation at every
  microsatellite in the genome — the MSI signature. The mutation rate
  at microsatellites can be $10^3$--$10^5$-fold elevated relative to
  non-repetitive sequences, even though the base substitution rate
  increases only $\sim$100-fold.

  \textbf{(c)} Probability that one cell acquires all $m$ mutations:
  \[
    p_\text{cell} \approx (\varepsilon \times n)^m.
  \]
  Probability that at least one cell in $N$ acquires all $m$:
  \[
    P \approx N \times p_\text{cell}.
  \]
  \begin{enumerate}[label=(\roman*)]
    \item MMR-proficient: $p_\text{cell} = (10^{-9}\times100)^5 = (10^{-7})^5 = 10^{-35}$;
          $P = 10^{10} \times 10^{-35} = 10^{-25}$. Essentially zero.
    \item MMR-deficient: $p_\text{cell} = (10^{-7}\times100)^5 = (10^{-5})^5 = 10^{-25}$;
          $P = 10^{10} \times 10^{-25} = 10^{-15}$. Still vanishingly small.
  \end{enumerate}

  \textbf{(d)} The simple model assumes the 5 mutations must occur
  \emph{independently in the same cell} from the same stem cell lineage.
  This is far too stringent. The crucial missing mechanism is
  \textbf{clonal expansion}: the first driver mutation confers a
  growth advantage, expanding the mutant clone to thousands or millions
  of cells. Each of these cells is now a candidate for acquiring the
  next mutation. This exponential amplification at each stage — the
  multi-stage Armitage--Doll model — collapses the time to accumulate
  all required mutations from thousands of years to decades.
  Lynch syndrome, by elevating $\varepsilon$ 100-fold, shortens each
  clonal expansion step, compressing cancer latency from $\sim$60--70
  years to $\sim$30--40 years — consistent with the observed earlier
  age of cancer onset in Lynch syndrome patients.
\end{solution}

% ============================================================
%  END OF MODULE 1
% ============================================================

\end{document}
